\documentclass[12pt,journal,onecolumn]{IEEEtran}
\usepackage{longtable}

\begin{document}

\title{Task 4: Prototype Implementation and Analysis}
\author{Team 27}
\maketitle

\section{Prototype Scope}

The prototype implements end-to-end non-trivial functionality for online interview preparation:

\begin{itemize}
    \item Role-based auth flow (student/admin)
    \item Timed test lifecycle (start, session fetch, sample run, violation, final submit)
    \item Async submission pipeline (API enqueue $\rightarrow$ worker evaluate $\rightarrow$ status update)
    \item Admin management for questions and users
    \item Student/admin analytics and monitoring endpoints
\end{itemize}

This satisfies the Task 4 requirement to demonstrate working architecture with significant end-to-end behavior.

\section{Implemented Architecture (Current)}

Current architecture pattern:
\begin{itemize}
    \item Modular monolith backend (Express + Mongoose modules)
    \item Event-driven evaluation path using BullMQ worker
    \item React frontend with role-protected routes
    \item Redis for queue + telemetry storage
    \item Docker sandbox for untrusted code execution
    \item Local MongoDB URI by default (\texttt{mongodb://localhost:27017/interview\_platform})
\end{itemize}

Key backend modules in code:
\begin{itemize}
    \item \texttt{auth}, \texttt{users}, \texttt{questions}, \texttt{tests}, \texttt{submissions}, \texttt{analytics}, \texttt{monitor}, \texttt{evaluation}, \texttt{shared}
\end{itemize}

\section{Performance and Reliability Tactics Used}

\subsection{Async queue decoupling}

Files:
\begin{itemize}
    \item \texttt{backend/src/submissions/submissions.routes.js}
    \item \texttt{backend/src/evaluation/queue.js}
    \item \texttt{backend/src/worker.js}
    \item \texttt{backend/src/evaluation/processSubmission.js}
\end{itemize}

Effect:
\begin{itemize}
    \item Submission API returns after persistence + enqueue (\texttt{202}), while heavy evaluation runs in worker.
\end{itemize}

\subsection{Cache-aside read path with invalidation}

Files:
\begin{itemize}
    \item \texttt{backend/src/shared/summaryCache.js}
    \item \texttt{backend/src/submissions/submissionReadService.js}
    \item \texttt{backend/src/questions/questions.routes.js}
\end{itemize}

Effect:
\begin{itemize}
    \item Hot GET APIs (submissions list, topics, summaries) reuse cached payloads.
    \item Write paths explicitly invalidate related cache keys.
\end{itemize}

\subsection{Startup warm-up for hot keys}

Files:
\begin{itemize}
    \item \texttt{backend/src/shared/cacheWarmup.js}
    \item \texttt{backend/src/server.js}
\end{itemize}

Effect:
\begin{itemize}
    \item Common read keys are loaded during startup to reduce cold-start read latency.
\end{itemize}

\subsection{Queue fallback for read reliability}

File:
\begin{itemize}
    \item \texttt{backend/src/submissions/submissionReadService.js}
\end{itemize}

Effect:
\begin{itemize}
    \item If queue telemetry is unavailable, submissions list falls back to DB-only response instead of failing endpoint.
\end{itemize}

\section{Architecture Comparison}

Compared architectures:

\begin{enumerate}
    \item Implemented: modular monolith + async worker queue
    \item Alternative: modular monolith + synchronous in-request evaluation (no queue/worker split)
\end{enumerate}

\begin{longtable}{|p{3.2cm}|p{5.3cm}|p{5.3cm}|}
\hline
\textbf{Aspect} & \textbf{Implemented (Async Queue)} & \textbf{Alternative (Synchronous API Evaluation)} \\
\hline
Submission request path & Returns after DB + enqueue & Request waits for full evaluation \\
\hline
API responsiveness under long code runs & Better, because evaluation is off request path & Worse, because request remains open until evaluator finishes/fails \\
\hline
Failure isolation & Worker failure isolated from API request thread & Evaluation failure directly impacts API response path \\
\hline
Recovery control & Queue attempts + dead-letter queue & Mainly client retry behavior \\
\hline
Operational complexity & Higher (Redis + worker lifecycle) & Lower (fewer moving parts) \\
\hline
\end{longtable}

Trade-off summary:
\begin{itemize}
    \item Implemented architecture improves responsiveness and isolation for heavy workloads.
    \item Alternative is simpler operationally but gives poorer user experience during long/failed evaluations.
\end{itemize}

\section{Quantified NFRs (System-Level Constraints)}

The following NFRs are written as measurable constraints with units and acceptance checks.

\subsection{NFR1: Performance}

System constraints:
\begin{itemize}
    \item Submission APIs return quickly because evaluation is asynchronous.
    \item Hard evaluator time budget per run is limited by \texttt{DOCKER\_TIMEOUT\_SEC = 10s}.
    \item Parallel evaluation capacity is \texttt{WORKER\_CONCURRENCY = 2}.
    \item Read latency control via cache: submissions list TTL \texttt{8000 ms}, topics TTL \texttt{60000 ms}.
\end{itemize}

Quantified targets:
\begin{itemize}
    \item API acknowledgement latency (\texttt{POST /api/submissions}, \texttt{POST /api/tests/\{id\}/submit}) p95 $\leq$ \texttt{500 ms} under normal local load (excluding DB/Redis outage).
    \item Queue wait estimate reported by API follows:
    \item \texttt{expectedQueueWaitSec = round(waiting * avgEvalMs / (workers * 1000))}
    \item Per question evaluation wall-time must be $\leq$ \texttt{10s} timeout budget (or marked failed/timeout).
\end{itemize}

Acceptance checks:
\begin{itemize}
    \item Verify p95 API latency from monitor snapshot (\texttt{/api/monitor/dashboard}).
    \item Verify completed submissions expose \texttt{actualProcessingSeconds}.
\end{itemize}

\subsection{NFR2: Security}

System constraints:
\begin{itemize}
    \item Untrusted code sandbox limits are fixed to CPU \texttt{1} core (\texttt{--cpus 1}), memory \texttt{256 MB} (\texttt{--memory 256m}), network disabled (\texttt{--network none}), timeout \texttt{10s}.
    \item JWT validity window default is \texttt{24h}.
    \item Public unauthenticated API surface is limited to \texttt{POST /api/auth/signup} and \texttt{POST /api/auth/signin}.
\end{itemize}

Quantified targets:
\begin{itemize}
    \item \texttt{100\%} of endpoints under \texttt{/api} except the two auth endpoints require valid JWT.
    \item \texttt{100\%} of code evaluations execute inside container with network disabled.
\end{itemize}

Acceptance checks:
\begin{itemize}
    \item Route audit in app middleware and auth routes.
    \item Docker command audit in evaluator runner for CPU/memory/network/timeout flags.
\end{itemize}

\subsection{NFR3: Reliability}

System constraints:
\begin{itemize}
    \item Queue retry policy: attempts per job \texttt{2}, retry backoff exponential with base delay \texttt{1000 ms}.
    \item Worker liveness telemetry: heartbeat interval \texttt{5000 ms}, stale threshold \texttt{30000 ms}.
\end{itemize}

Quantified targets:
\begin{itemize}
    \item Transient queue/worker failures should not cause immediate data loss because failed jobs get one retry window.
    \item Read path degrades gracefully: if queue telemetry fails, submissions list still returns DB-backed response.
\end{itemize}

Acceptance checks:
\begin{itemize}
    \item Simulate queue telemetry failure and verify submissions API still responds successfully.
    \item Verify failed jobs increment attempts and retry with configured backoff.
\end{itemize}

\subsection{NFR4: Availability}

System constraints:
\begin{itemize}
    \item Readiness endpoint is strict dependency gate: \texttt{/health/ready} returns \texttt{503} when Mongo/Redis/worker readiness is not satisfied.
    \item Liveness endpoint remains a process health check.
\end{itemize}

Quantified targets:
\begin{itemize}
    \item Readiness accuracy target: \texttt{100\%} dependency-aware status reporting.
    \item Operational target: monthly availability $\geq$ \texttt{99.5\%} when dependencies are healthy.
\end{itemize}

Acceptance checks:
\begin{itemize}
    \item Stop Redis or worker and verify \texttt{/health/ready} returns \texttt{503}.
    \item Restore dependencies and verify \texttt{/health/ready} returns \texttt{200}.
\end{itemize}

\section{Evidence Map (Key Files)}

\begin{itemize}
    \item Auth gate and health: \texttt{backend/src/app.js}
    \item Queue config: \texttt{backend/src/evaluation/queue.js}
    \item Worker lifecycle: \texttt{backend/src/worker.js}
    \item Evaluation strategies/factory: \texttt{backend/src/evaluation/strategies/*}, \texttt{backend/src/evaluation/strategyFactory.js}
    \item Docker sandbox runner: \texttt{backend/src/evaluation/dockerRunner.js}
    \item Cache core: \texttt{backend/src/shared/summaryCache.js}
    \item Cache warm-up: \texttt{backend/src/shared/cacheWarmup.js}
    \item Submission read path: \texttt{backend/src/submissions/submissionReadService.js}
    \item Topics cache + invalidation: \texttt{backend/src/questions/questions.routes.js}
    \item Monitor endpoint: \texttt{backend/src/monitor/monitor.routes.js}
\end{itemize}

\section{Conclusion}

Task 4 prototype architecture is aligned with the implemented codebase and supports the core non-trivial workflow.

Compared with synchronous in-request evaluation, the current async queue architecture gives stronger response-time behavior and better failure isolation, while introducing manageable infrastructure complexity.

\end{document}
